%! TeX root = ../phd-thesis.tex
%! Author = davidedomini

%----------------------------------------------------------------------------------------
\chapter{Federated Learning}
\label{chap:federated-learning}
\minitoc
%----------------------------------------------------------------------------------------

\gls{FL} is a distributed machine-learning paradigm 
 in which multiple clients collaboratively train one or more models 
 while retaining their training data locally. 
% 
Rather than collecting observations in a centralized dataset, 
 \gls{FL} moves the learning procedure to the locations where data are generated 
 and exchanges model parameters or updates during 
 training~\cite{mcmahan2017communication,kairouz2021advances}. 
% 
This organization is particularly relevant when data are private, 
costly to transfer, continuously generated at the edge, 
or distributed across independently managed devices and organizations.

The distribution of the learning process introduces challenges that do not arise,
 or are less prominent, in conventional centralized training. 
% 
The system must coordinate multiple local optimization procedures, 
 communicate model updates, handle clients that may participate intermittently, 
 and learn from datasets whose statistical properties can differ substantially. 
% 
Moreover, 
 although the original formulation of \gls{FL} assumes a central coordinating server, 
  the separation between data ownership and model training does not prescribe 
  a particular communication architecture. 
%  
Federated optimization can also be organized through hierarchical,
 peer-to-peer, or fully decentralized interactions.



\section{Training Dynamics}
\label{sec:fl-training-dynamics}

\glsentryfull{FL} organizes model training across a set of \(K\) clients,
 each of which retains exclusive access to a local dataset \(\mathcal{D}_k\). 
% 
The size of the local dataset and the total number of training samples are respectively defined as
%
\begin{equation}
    n_k = \lvert\mathcal{D}_k\rvert,
    \qquad
    n = \sum_{k=1}^{K} n_k.
    \label{eq:fl-dataset-sizes}
\end{equation}
%
Given a model parameterized by \(\boldsymbol{\theta}\) 
 and a sample-wise loss function \(\ell\), 
 client \(k\) defines the local empirical objective
%
\begin{equation}
    F_k(\boldsymbol{\theta})
    =
    \frac{1}{n_k}
    \sum_{(\mathbf{x},y)\in\mathcal{D}_k}
    \ell(\boldsymbol{\theta};\mathbf{x},y).
    \label{eq:fl-local-objective}
\end{equation}
%
In the simplest \gls{FL} formulation, 
 all clients collaborate to learn a single global model 
 by minimizing the weighted combination of their local objectives:
%
\begin{equation}
    \boldsymbol{\theta}^{*}
    =
    \operatorname*{arg\,min}_{\boldsymbol{\theta}}
    F(\boldsymbol{\theta}),
    \qquad
    F(\boldsymbol{\theta})
    =
    \sum_{k=1}^{K}
    \frac{n_k}{n}
    F_k(\boldsymbol{\theta}).
    \label{eq:fl-global-objective}
\end{equation}
%
The weights account for differences in the sizes of the local datasets. 
%
The objective in~\Cref{eq:fl-global-objective} corresponds 
 to empirical risk minimization over all available samples,
  but its optimization is decomposed into computations performed directly by the clients, 
  without requiring their datasets to be transferred to a common location.

The basic realization of this process adopts a client-server architecture 
 and the \gls{FedAvg} algorithm~\cite{mcmahan2017communication}. 
% 
A central server coordinates training and maintains the current global parameters,
 while the clients perform local optimization. 
% 
Training proceeds through a sequence of communication rounds.
%
At round \(t\), 
 the server holds the global parameters \(\boldsymbol{\theta}^{t}\) 
 and performs the following operations:

\begin{enumerate}
    \item \emph{Client selection.}
    The server selects a subset \(\mathcal{S}_{t}\) of the available clients.
    Partial participation is common when the client population is large or
    individual clients are only intermittently available.

    \item \emph{Model distribution.}
    The server sends \(\boldsymbol{\theta}^{t}\) to every selected client, so
    that all participants start the round from the same model state.

    \item \emph{Local optimization.}
    Each selected client independently trains the received model on its local
    dataset. For example, using mini-batch stochastic gradient descent, client
    \(k\) repeatedly applies
    %
    \begin{equation}
        \boldsymbol{\theta}_{k}^{t,e+1}
        =
        \boldsymbol{\theta}_{k}^{t,e}
        -
        \eta
        \nabla
        \ell
        \left(
            \boldsymbol{\theta}_{k}^{t,e};
            \mathcal{B}_{k}^{e}
        \right),
        \label{eq:fl-local-update}
    \end{equation}
    %
    where \(e\) identifies a local optimization step, \(\eta\) is the learning
    rate, and
    \(\mathcal{B}_{k}^{e}\subseteq\mathcal{D}_{k}\) is a local mini-batch.

    \item \emph{Update collection.}
    After completing one or more local epochs, each selected client sends the
    resulting model parameters, or their difference from the received global
    parameters, back to the server.

    \item \emph{Aggregation.}
    The server combines the received models to obtain the global parameters for
    the next round. FedAvg computes a weighted average:
    %
    \begin{equation}
        \boldsymbol{\theta}^{t+1}
        =
        \sum_{k\in\mathcal{S}_{t}}
        \frac{n_k}{
            \sum_{j\in\mathcal{S}_{t}} n_j
        }
        \boldsymbol{\theta}_{k}^{t+1}.
        \label{eq:fedavg-aggregation}
    \end{equation}
\end{enumerate}

\begin{figure}[t]
    \centering
    \fbox{
        \parbox[c][5.2cm][c]{0.9\textwidth}{
            \centering
            \textbf{Image placeholder: client--server FL training round}
            \\[1em]
            1. The server selects participating clients and distributes the
            current global model.
            \\[0.5em]
            2. Clients independently train the model on their local datasets.
            \\[0.5em]
            3. Clients return their locally updated models.
            \\[0.5em]
            4. The server aggregates the updates and starts a new round.
        }
    }
    \caption{Training dynamics of client--server Federated Learning. Raw data
    remain on the clients, while model parameters or updates are exchanged
    during successive communication rounds.}
    \label{fig:fl-training-round}
\end{figure}

The aggregated parameters become the initial model for the following round. 
%
The process continues until a stopping condition is reached, 
 such as a fixed number of rounds, stabilization of a validation metric, 
 or exhaustion of the available training budget. 
% 
Local computation and global communication are therefore interleaved 
 throughout the entire learning process.

The number of local optimization steps controls an important trade-off.
%
Performing more local work between two communication rounds reduces 
 the frequency with which models must be exchanged, 
 but allows each local model to move farther from the shared parameters. 
% 
This effect is usually limited when clients observe similar data distributions, 
 whereas it can become problematic when their local objectives differ. 
% 
The number of participating clients, their dataset sizes, local learning rates, 
 and aggregation weights similarly affect the relationship between 
 the local optimization processes and the global objective.

Federated settings are commonly distinguished according to the scale 
 and availability of their clients. 
% 
In \emph{cross-silo} \gls{FL}, 
 the participants are typically a limited number of stable organizations, 
 such as hospitals or companies, with relatively large datasets and reliable infrastructure. 
% 
In \emph{cross-device} \gls{FL}, 
 training may involve a much larger population of phones, sensors, vehicles, 
 or edge devices that hold smaller datasets and may be available only intermittently. 
% 
The optimization principle remains the same, 
 but client selection, fault tolerance, communication cost, and resource variability
 become substantially more important in the cross-device case~\cite{kairouz2021advances}.

\section{Decentralized Architectures}
\label{sec:fl-decentralized}

\gls{FL} and centralized coordination are distinct concepts. 
%
The term \emph{federated} describes how data ownership and training are organized:
 participants retain their datasets and cooperate 
 by exchanging learning-related information. 
%
It does not require a central server. 
%
The client-server architecture is widely adopted because it provides 
 a simple way to initialize the model, coordinate rounds, collect updates,
 and construct a consistent global state, but the same learning objective can be realized 
 through different communication topologies.

The central server introduces an explicit coordination point. 
%
This simplifies model versioning and aggregation, 
 since every round has a clearly identified initial and final state. 
% 
However, 
 it can also become a communication bottleneck when many clients participate, 
 and its unavailability may suspend the entire training process. 
% 
Moreover, 
 all clients must be able to communicate with the server even when 
 direct local communication would be less expensive or more readily available.

Decentralized \gls{FL} instead represents the clients 
 and their communication relationships as a graph
%
\begin{equation}
    G = (V,E),
\end{equation}
%
where each vertex \(i\in V\) is a client and an edge \((i,j)\in E\) indicates
that clients \(i\) and \(j\) can exchange model information. A client typically
alternates local optimization with a mixing operation over its neighborhood
\(\mathcal{N}_{i}\). A generic mixing step can be written as
%
\begin{equation}
    \boldsymbol{\theta}_{i}^{t+1}
    =
    \sum_{j\in\mathcal{N}_{i}\cup\{i\}}
    W_{ij}
    \widehat{\boldsymbol{\theta}}_{j}^{t},
\end{equation}
%
where \(\widehat{\boldsymbol{\theta}}_{j}^{t}\) is the model obtained by client
\(j\) after local training and \(W_{ij}\) is the weight assigned by \(i\) to
that model. The matrix \(\mathbf{W}\) encodes the communication topology and
aggregation policy. Under appropriate connectivity and weighting assumptions,
repeated local exchanges allow information to propagate through the network and
can drive the clients toward a consensus model.

\begin{figure}[t]
    \centering
    \fbox{
        \parbox[c][5.4cm][c]{0.9\textwidth}{
            \centering
            \textbf{Image placeholder: federated-learning architectures}
            \\[1em]
            Client--server: clients communicate with one global aggregator.
            \\[0.5em]
            Hierarchical: clients communicate with intermediate edge aggregators.
            \\[0.5em]
            Peer-to-peer: clients exchange models directly over a communication
            graph.
            \\[0.5em]
            Gossip-based: clients repeatedly exchange and mix models with a
            limited set of neighbors.
        }
    }
    \caption{Representative runtime architectures for Federated Learning. The
    organization of communication and aggregation can vary without changing the
    fundamental constraint that training data remain distributed among clients.}
    \label{fig:fl-architectures}
\end{figure}

Several decentralized organizations can be derived from this model. In a fully
connected peer-to-peer topology, every client can exchange its update with every
other client and reproduce an aggregation similar to the server-based case.
Although conceptually simple, all-to-all communication becomes expensive as the
population grows. Sparse peer-to-peer systems restrict exchanges to logical or
physical neighbors and propagate information over multiple steps
\cite{wink2021approach}. Gossip-based approaches further reduce coordination
requirements by letting clients periodically communicate with one or a few
peers, possibly without a globally synchronized round structure
\cite{hegedus2021decentralized}.

Hierarchical architectures introduce intermediate aggregators, commonly located
on edge or regional infrastructure. Clients first send updates to a nearby
aggregator; intermediate models are subsequently combined at higher levels.
This organization can reduce long-distance communication and distribute the
aggregation workload, while still retaining some designated coordination
entities. Hybrid architectures may combine server-based coordination in stable
parts of the network with direct exchanges among devices that cannot
continuously reach the central infrastructure.

Removing the global server distributes its responsibilities rather than
eliminating them. Clients must establish which models can be combined, manage
possibly stale model versions, and determine when sufficient information has
been received. Convergence then depends not only on the optimization algorithm,
but also on graph connectivity, communication frequency, message delays, and
node churn. Failures may have a more localized effect than in the client--server
case, but temporary partitions can prevent information from spreading between
different parts of the network.

The architecture also determines the meaning of a training round. In a
synchronous peer-to-peer system, all clients may alternate between a common
local-training phase and a common communication phase. In asynchronous and
gossip-based systems, each client instead progresses according to its own local
clock and aggregates models produced at different times. The latter organization
avoids waiting for the slowest clients but requires mechanisms for handling
staleness and preventing rapidly communicating clients from having a
disproportionate influence.

Importantly, decentralization alone does not solve statistical heterogeneity.
A peer-to-peer network can still drive every client toward a single consensus
model, reproducing the same statistical compromise as centralized aggregation.
Conversely, multiple specialized models can be maintained by a centralized
server. Communication architecture determines how coordination is performed;
the number and specialization of the learned models determine what learning
objective the coordination realizes.


\section{Data Heterogeneity}
\label{sec:fl-heterogeneity}

The global objective introduced in
Section~\ref{sec:fl-training-dynamics} does not require local datasets to have
the same size, but its simplest behavior is obtained when all observations are
sampled independently from the same underlying probability distribution. In an
IID setting,
%
\[
    (\mathbf{x},y)\sim P
\]
%
for every client, so differences between local datasets are mainly due to finite
sampling. Local gradients then provide noisy but broadly aligned estimates of
the gradient of the global objective.

Real federated datasets rarely satisfy this assumption. Each client \(k\) may
instead observe a different distribution
%
\[
    (\mathbf{x},y)\sim P_k(\mathbf{x},y),
\]
%
with \(P_i\neq P_j\) for at least some pair of clients \(i\) and \(j\). This
condition is generally referred to as \emph{non-independent and identically
distributed}, or \emph{non-IID}, data. The term covers several distinct forms of
heterogeneity \cite{zhu2021federated}:

\begin{itemize}
    \item \emph{Label-distribution skew} occurs when the frequency of the target
    classes differs among clients. Some clients may observe only a subset of the
    possible labels, while others may contain the same labels in different
    proportions.

    \item \emph{Feature-distribution skew} occurs when clients share the same
    prediction task but observe different input distributions. Images may, for
    example, differ in lighting, background, sensor characteristics, or writing
    style even when their labels are drawn from the same set.

    \item \emph{Concept shift} occurs when the relationship between inputs and
    targets differs among clients. Formally, conditional distributions such as
    \(P_k(y\mid\mathbf{x})\) are client-dependent, so the same input may require
    a different prediction in different contexts.

    \item \emph{Quantity skew} occurs when clients hold substantially different
    numbers of observations. This condition does not necessarily change the
    underlying distributions, but it affects aggregation weights, uncertainty,
    and the influence of individual clients.
\end{itemize}

Multiple forms can coexist. For example, a client can hold fewer observations,
cover only a subset of the labels, and acquire its inputs through a sensor whose
characteristics differ from those of the remaining devices. Non-IID behavior
should therefore not be treated as a single scalar property of a dataset.

Statistical heterogeneity affects federated optimization because the local
objectives \(F_k\) no longer have the same minimizers. A client that performs
several local steps moves its parameters toward a solution that is effective for
its own distribution, which may differ from the direction preferred by other
clients. This phenomenon is commonly referred to as \emph{client drift}. When
the resulting updates are averaged, they may partially cancel one another,
produce unstable progress, or converge to a compromise that performs unevenly
across the client population \cite{li2022federated}. Increasing the amount of
local training can amplify this divergence even though it reduces the number of
communication rounds.

Non-IID data may be unstructured, with each client receiving an unrelated local
distribution, but it can also exhibit a regular organization. This is especially
important in spatially distributed systems. Consider an environment
\(\mathcal{A}\) partitioned into \(Q\) regions,
%
\[
    \mathcal{A}
    =
    \bigcup_{q=1}^{Q} a_q,
\]
%
where each region \(a_q\) is associated with a data-generating distribution
\(P_q\). A client located in \(a_q\) observes data sampled primarily from
\(P_q\). Clients in the same region therefore tend to have statistically similar
datasets, whereas clients located in different regions may observe different
distributions:
%
\[
    P_i \simeq P_j
    \quad
    \text{if clients \(i\) and \(j\) are in the same region},
    \qquad
    P_i \not\simeq P_j
    \quad
    \text{across heterogeneous regions}.
\]

\begin{figure}[t]
    \centering
    \fbox{
        \parbox[c][5.5cm][c]{0.9\textwidth}{
            \centering
            \textbf{Image placeholder: spatially structured non-IID data}
            \\[1em]
            Represent an environment divided into colored regions, with each
            color denoting a different data distribution.
            \\[0.6em]
            Place multiple devices inside each region to show that nearby
            devices tend to observe similar distributions, while devices in
            different regions observe non-IID data.
        }
    }
    \caption{Spatially structured data heterogeneity. Data are comparatively
    homogeneous within each region but heterogeneous across regions, producing
    group-level rather than purely client-level non-IID behavior.}
    \label{fig:spatial-non-iid}
\end{figure}

This organization can arise whenever observations depend on the surrounding
environment. Traffic measurements vary between residential areas, city centers,
and highways; energy consumption depends on building type and local climate;
air-quality measurements are affected by nearby infrastructure; and images
collected by mobile devices depend on the locations and conditions in which they
operate. Spatial proximity is therefore often correlated with statistical
similarity, although it is not a guarantee of it. Nearby devices may perform
different tasks or face different orientations and sensing conditions, while
distant areas may still generate similar data.

Spatial heterogeneity can also change over time. A mobile client may move
between regions and progressively replace observations from one distribution
with observations from another. Even stationary devices may experience temporal
changes caused by weather, recurring activity patterns, infrastructure changes,
or other environmental processes. The local distribution is then more accurately
represented as \(P_k^t\), making heterogeneity both spatial and temporal.

Several families of FL algorithms attempt to mitigate these effects. Optimization
methods can constrain local updates or correct the disagreement between local
and global directions, as done respectively by FedProx and SCAFFOLD
\cite{li2020federated,karimireddy2020scaffold}. Other approaches personalize a
shared model for each client, learn common representations with local output
components, or allow different groups of clients to train different models.
Clustered Federated Learning follows the last strategy and exploits structure in
the client population rather than forcing every distribution into a single global
objective.


\section{Clustered Federated Learning}
\label{sec:clustered-fl}

Clustered Federated Learning (CFL) is a form of personalized FL in which clients
are partitioned into groups and a separate model is trained for each group
\cite{tan2023towards}. It occupies an intermediate position between conventional
FL, which trains one model for the entire population, and purely local learning,
which trains an independent model for every client. Clients within a cluster
continue to benefit from collaborative training, while clients associated with
substantially different distributions are not forced to share the same final
parameters.

Let
%
\[
    \mathcal{C}
    =
    \left\{
        C_1,C_2,\ldots,C_Q
    \right\}
\]
%
be a partition of the \(K\) clients, where
\(C_q\cap C_r=\emptyset\) for \(q\neq r\) and
\(\bigcup_{q=1}^{Q}C_q=\{1,\ldots,K\}\). Each cluster \(C_q\) is assigned a
model \(\boldsymbol{\theta}_q\) obtained by minimizing its cluster-specific
objective:
%
\[
    \boldsymbol{\theta}_q^{*}
    =
    \operatorname*{arg\,min}_{\boldsymbol{\theta}_q}
    F_q(\boldsymbol{\theta}_q),
\]
%
where
%
\[
    F_q(\boldsymbol{\theta}_q)
    =
    \sum_{k\in C_q}
    \frac{n_k}{N_q}
    F_k(\boldsymbol{\theta}_q),
    \qquad
    N_q
    =
    \sum_{k\in C_q}n_k.
\]
%
The intended partition places clients with compatible local objectives in the
same cluster. Exact IID behavior within a cluster is neither necessary nor
generally attainable; the practical goal is to make intra-cluster heterogeneity
smaller than the heterogeneity of the complete population.

A basic clustered training process alternates between determining which clients
should collaborate and optimizing the corresponding cluster models. It can be
organized through the following phases:

\begin{enumerate}
    \item \emph{Model initialization.}
    One or more candidate models are initialized. They may start from common
    parameters or from independently initialized states.

    \item \emph{Client assignment.}
    Each client is associated with the model or group considered most compatible
    with its local data.

    \item \emph{Local training.}
    Clients train the model associated with their cluster using their local
    datasets.

    \item \emph{Cluster-wise aggregation.}
    Updates produced by members of the same cluster are aggregated to obtain a
    new cluster model. Updates belonging to different clusters remain separate.

    \item \emph{Cluster revision.}
    Client assignments can be reevaluated, and clusters may be split, merged, or
    otherwise reorganized as additional evidence about their members becomes
    available.
\end{enumerate}

Some methods begin with a common model and split the population when client
updates become sufficiently dissimilar
\cite{sattler2021clustered}. Others maintain a predefined set of candidate models
and alternate between assigning each client to its best-performing model and
updating the models from the resulting assignments
\cite{ghosh2022efficient}. In the latter case, a client can select its cluster
according to the local empirical loss:
%
\[
    c_k^{t}
    =
    \operatorname*{arg\,min}_{q\in\{1,\ldots,Q\}}
    F_k(\boldsymbol{\theta}_q^{t}).
\]
%
The selected model is then trained locally, and its update contributes only to
the aggregation for cluster \(c_k^t\).

Client similarity can be estimated in several ways. Loss-based strategies
evaluate candidate cluster models on local data and assign the client to the
model with the lowest loss. Gradient-based strategies compare the direction or
magnitude of local updates, under the assumption that compatible objectives
produce aligned optimization trajectories. Weight-based strategies compare
model parameters or compact representations derived from them. Additional
metadata, such as location, communication reachability, device role, or sensing
context, can be used as a prior when it is expected to correlate with the data
distribution.

CFL is particularly interesting when heterogeneity has a group-level structure.
If several clients observe different samples from the same underlying
distribution, training them independently would discard useful opportunities for
data sharing. Conversely, aggregating their updates with those of clients from
unrelated distributions may degrade their models. Clustering preserves
collaboration within statistically coherent groups while allowing the resulting
models to specialize. In spatial settings, the clusters may correspond to
geographic zones or other environmental contexts, although client location and
statistical similarity should remain conceptually distinct.

Clustered training can be implemented using either centralized or decentralized
architectures. A central server may maintain all cluster models, evaluate
membership information, and perform separate aggregations. In a decentralized
implementation, clients can discover compatible peers and exchange updates only
within the resulting groups. Hierarchical systems may assign one regional
aggregator to each cluster. Consequently, CFL specifies which models should be
learned and which clients should contribute to them, whereas the architecture
specifies how those clients coordinate.

The approach also introduces challenges absent from single-model FL. The number
of clusters may be unknown, and an incorrect value can either merge incompatible
distributions or fragment compatible clients into unnecessarily small groups.
Similarity estimates obtained early in training can be noisy because models have
not yet learned sufficiently informative representations. Repeated reassignment
can cause unstable memberships, while fixed assignments cannot accommodate
clients whose distributions change over time.

Clustering also has computational and communication costs. Multiple models must
be initialized, stored, evaluated, and distributed; update comparisons may reveal
additional information about local training behavior; and small clusters may not
contain enough data to train a reliable specialized model. Hard assignments force
each client to select one model even when its data combine multiple distributions,
whereas soft or overlapping assignments require more complex aggregation.
Mobility and temporal drift further require the system to decide whether a client
should change clusters, adapt its current model, or retain knowledge acquired in
previous environments.

CFL therefore does not eliminate data heterogeneity. Instead, it changes the
learning objective so that structured heterogeneity can be represented by
multiple collaborative models. Its effectiveness depends on whether the client
population contains meaningful groups, whether those groups can be identified
from information available during federated training, and whether the clustering
process can remain stable as the system evolves.